\documentclass[10pt,a4paper]{amsart}
\usepackage[margin=19mm]{geometry}
\usepackage{amsmath,amssymb,mathtools}
\usepackage[scr=rsfs,cal=euler]{mathalfa}
\usepackage{xfrac}
\usepackage[unicode,hidelinks,bookmarks=false]{hyperref}
\newcommand{\ie}{i.e.\ }
\newcommand{\cf}{cf.\ }
\newcommand{\resp}{respectively\ }
\newcommand{\Subsection}{\subsection}
\providecommand{\nobreakdash}{\nobreak\hbox{-}}
\setlength{\emergencystretch}{2em}
\multlinegap0pt
\usepackage[shortlabels]{enumitem}
\usepackage{amssymb}
\usepackage{xcolor}
\usepackage{bbm}
\usepackage{enumerate}
\usepackage{tikz}
\newtheorem{lem}{Lemma}
\newcommand{\R}{\mathbb{R}} 
\newcommand{\Z}{\mathbb{Z}}
\title{Extremal diameters of 3-coloring graphs of trees}
\author{Shamil Asgarli, Sara Krehbiel, Simon MacLean, Gjergji Zaimi}
\date{Source: arXiv:2512.03789v2 (CC BY 4.0)}
\begin{document}
\maketitle

\noindent\emph{Source and licence.} Extremal diameters of 3-coloring graphs of trees. Version arXiv:2512.03789v2 (CC BY 4.0), distributed by arXiv under the Creative Commons Attribution 4.0 International licence, http://creativecommons.org/licenses/by/4.0/, as recorded in the arXiv licence metadata for this exact version and shown on the abstract page https://arxiv.org/abs/2512.03789v2. Full licence notice: \texttt{licenses/arxiv-2512.03789-CC-BY-4.0.txt}.\par\smallskip

\noindent\textit{Editorial scope.} Counted unit: the article's lem lem:median-control (source0.tex lines 346--348). The article's complete original proof of it is delivered with it (source0.tex lines 350--405, one \texttt{proof} environment of 56 lines). No mathematics is written, repaired or re-derived: the statement and the proof are the article's own lines, in the article's own notation, and no retained line is substituted or re-flowed.  Retained uncounted as the source-local context the proof needs: lines 321--324; lines 335--338.  Nothing was omitted from any retained span, so the package carries no omission record.  No retained line contains a citation, so the wrapper carries no bibliography and no bibliography entry.  No retained line contains a figure environment, an image-inclusion command or drawing code, so no image is delivered and the package declares no asset dependency.  Editorial formatting only, in addition: an AMS \texttt{amsart} preamble that restates the article's own declarations and macros used by the retained lines, namely the environments \texttt{lem} on separate counters, together with the standard packages those lines need.  The retained environments print in this document as Lemma 1 in order of appearance, whereas the article prints the same items with the numbering of its own full document.  The complete native source of the article as distributed in the arXiv e-print for this exact version is bundled unchanged as \texttt{sources/190285/source0.tex}.
\begin{lem}[Median and minimizer]
\label{lem:minimizer}
Let $M$ be the median of the labels of $h$. An integer $k \in \Z$ minimizes $\Phi(k)$ if and only if $-k\in [M-\frac{1}{2}, M+\frac{1}{2}]$. In particular, if $|M| < 1$, then $0$ is a global minimizer of $\Phi$.
\end{lem}

\begin{lem}
\label{lem:lattice}
Let $f\colon \Z \to \R$ be convex with a global minimum at $0$. Let $L \subset \Z$ be a lattice (i.e., an arithmetic progression) such that $0 \notin L$, containing points strictly less than and strictly greater than $0$. Then $\min_{x \in L} f(x)$ is attained at one of the two lattice points $u < 0 < v$ of $L$ adjacent to $0$.
\end{lem}

\begin{lem}\label{lem:median-control}
Suppose $h$ is a balanced labeling of $T$ with maximum $L_1$-norm (among all balanced labelings of $T$). Let $M$ be the median of the labels of $h$. Then $M\in \{1, \frac{3}{2}, 2\}$.
\end{lem}

\begin{proof} By symmetry (replacing $h$ with $-h$ if necessary), we assume $M \ge 0$. Recall that $a_i$ denotes the number of vertices with label $i$.

\medskip\noindent
\textbf{Upper Bound ($M \le 2$).}
Since $h$ is balanced, we have $\Phi(0) \le \Phi(-3)$. We expand the difference 
\[
\Phi(0) - \Phi(-3) = \sum_{i} a_i (|i| - |i-3|)\leq 0
\]
using:
\begin{itemize}[noitemsep]
    \item $i \le 0 \implies |i| - |i-3| = -3$; \quad $i \ge 3 \implies |i| - |i-3| = 3$.
    \item $i = 1 \implies |i| - |i-3| = -1$; \quad $i = 2 \implies |i| - |i-3| = 1$;
\end{itemize}
Thus, $-3A_{\le 0} - a_1 + a_2 + 3A_{\ge 3} \le 0$. Substituting $A_{\le 0} = n - a_1 - a_2 - A_{\ge 3}$ yields:
\[
-3n + 2a_1 + 4a_2 + 6A_{\ge 3} \le 0 \implies 6A_{\ge 3} \le 3n - 2a_1 - 4a_2 \le 3n.
\]
This implies $A_{\ge 3} \le n/2$. If $M > 2$, then $M \ge 3$ and so $A_{\ge 3} > n/2$, a contradiction.

\medskip\noindent
\textbf{Lower Bound ($M \ge 1$).}
Suppose $|M| < 1$. By Lemma \ref{lem:minimizer}, we know that $x=0$ is a global minimizer of $\Phi$, so slopes satisfy $\Delta_{-1} \le 0 \le \Delta_0$. By maximality, if $\|h+k\|_1 > \|h\|_1$, then $h+k$ is unbalanced.

\smallskip
\textbf{Case 1:} Strict minimum ($\Delta_{-1} < 0 < \Delta_0$).
Since $\|h\pm 1\|_1 > \|h\|_1$, both $h+1$ and $h-1$ are unbalanced.
For $h+1$ (on lattice $3\Z+1$ straddling $-2$ and $1$):
\[
\text{$h+1$ unbalanced} \implies \min(\Phi(-2), \Phi(1)) < \Phi(1) \implies \Phi(-2) < \Phi(1).
\]
For $h-1$ (on lattice $3\Z-1$ straddling $-1$ and $2$):
\[
\text{$h-1$ unbalanced} \implies \min(\Phi(-1), \Phi(2)) < \Phi(-1) \implies \Phi(2) < \Phi(-1).
\]
Summing these gives $\Phi(1) - \Phi(2) > \Phi(-2) - \Phi(-1)$, or $\Delta_1 < \Delta_{-2}$, contradicting convexity.

\smallskip
\textbf{Case 2:} Flat minimum (at least one slope is 0).
By symmetry ($h \mapsto -h$), we may assume $\Delta_{-1} = 0 \le \Delta_0$. This implies $A_{\ge 1} = (\Delta_{-1}+n)/2 = n/2$.
\begin{itemize}
    \item If $\Delta_0 = 0$, then $A_{\ge 0} = n/2$, so $a_0 = 0$. Since $A_{\ge 1} = n/2 > 0$ and $A_{\le -1} = n-A_{\geq 1} > 0$, the label set is \emph{not} contiguous, a contradiction.
    
    \item If $\Delta_0 > 0$, then $a_0 > 0$ and $a_1 > 0$. We compute the relevant slopes:
    \begin{align*}
        \Delta_0 &= 2A_{\ge 0} - n = 2(A_{\ge 1} + a_0) - n = 2a_0 > 0, \\ 
        \Delta_{-2} &= 2A_{\ge 2} - n = 2(A_{\ge 1} - a_1) - n = -2a_1 < 0.
    \end{align*}
    Thus, $\|h+1\|_1 > \|h\|_1$ and $\|h-2\|_1 > \|h\|_1$. By maximality, both are unbalanced on the lattice $3\Z+1$. Applying Lemma~\ref{lem:lattice}, we obtain:
    \[
       \Phi(-2) < \Phi(1) \quad (\text{from } h+1) 
       \quad \text{and} \quad
       \Phi(1) < \Phi(-2) \quad (\text{from } h-2),
    \]
    a contradiction.
\end{itemize}
Thus, the lower bound $|M| \ge 1$ holds. \end{proof}

\end{document}
